% \codeavail is defined in main.tex (named build: the repository URL; review build: anonymized).
% The fallback below only keeps this file compiling if main.tex does not define it.
\providecommand{\codeavail}{available at \expandafter\url\expandafter{\repourl}}

\section{Verification and Reproducibility}\label{sec:verify}

\emph{Cross-checks.} Our CGT engine agrees with a separate brute-force sum solver on $55{,}195$ positions $G$
(at most two pawns a side, $3$--$5$ files), for three race rules and $15$ reference games $R$: all
$4{,}967{,}550$ outcomes of $G+R$ match under the textbook test, e.p.\ on and off, and all $2{,}483{,}775$
under the e.p.-aware test. A whole-board minimax confirms its outcomes on $1{,}070$ single components. A Lean~4
project checks $123$ finite statements (Section~S17). AI models wrote the full \lab{S} proofs. AI referees, some from another vendor, checked them, and
programs ran the ring base cases; no human referee has yet.

\emph{Code and data availability.} The code and data are \codeavail. The supplementary material gives the
full proofs of the \lab{S} results (Sections~S12--S15) and the FEN of the main positions with kings (Section~S11).
